%% toc: P1 2 2.1
\section{Proofs for membership, set operations, and bounded set formation}
\label{sec:annex-p01-w02-p01-005-set-operations}

Throughout, ``by extensionality'' refers to clause 3 of
\cref{def:p01-w02-p01-005-membership}: two sets are equal iff they have the
same members. The variables range over sets, and the ambient logic is
classical, so for every set \(z\) and every set \(A\) exactly one of
\(z\in A\) and \(z\notin A\) holds.

\begin{annexitem}{annex:p01-w02-p01-005-separated-set}{thm:p01-w02-p01-005-separated-set}{Existence, uniqueness and inclusion of a separated set}
\begin{proof}
\emph{Existence.} By the side conditions, \(\varphi(x,\bar p)\) is a formula of
\(\Lang_{\in}\) with free variables among \(x,\bar p\), and the variable \(b\)
used for the formed set does not occur free in it. The separation schema
therefore has the instance
\[
  \forall a\,\exists b\,\forall x\,
  \bigl(x\in b\leftrightarrow(x\in a\land\varphi(x,\bar p))\bigr)
\]
for the assigned parameter sets \(\bar p\). Instantiating \(a\) by the set
\(A\) gives a set \(b_0\) with
\(\forall x\,(x\in b_0\leftrightarrow(x\in A\land\varphi(x,\bar p)))\).

\emph{Uniqueness.} Let \(b_1\) be a set with the same property. For an
arbitrary set \(x\), the two properties give
\(x\in b_0\leftrightarrow(x\in A\land\varphi(x,\bar p))\) and
\((x\in A\land\varphi(x,\bar p))\leftrightarrow x\in b_1\), hence
\(x\in b_0\leftrightarrow x\in b_1\). Since \(x\) was arbitrary,
\(b_0\) and \(b_1\) have the same members, and \(b_0=b_1\) by extensionality.

\emph{Inclusion.} For an arbitrary set \(x\), if \(x\in b_0\) then
\(x\in A\land\varphi(x,\bar p)\), and in particular \(x\in A\). Thus
\(\forall x\,(x\in b_0\to x\in A)\), which is \(b_0\subseteq A\) by the
definition of inclusion in \cref{sec:set-prerequisites-01}.

\emph{Renaming.} Let \(y\) not occur in \(\varphi(x,\bar p)\), and let
\(\varphi(y,\bar p)\) result from replacing the free occurrences of \(x\) by
\(y\). The formula \(\varphi(y,\bar p)\) has \(b\) not free, so by existence
and uniqueness \(\{y\in A:\varphi(y,\bar p)\}\) is the unique set \(b_2\) with
\(\forall y\,(y\in b_2\leftrightarrow(y\in A\land\varphi(y,\bar p)))\). Fix
a set \(z\). Interpreting \(y\) by \(z\) in \(\varphi(y,\bar p)\) and
interpreting \(x\) by \(z\) in \(\varphi(x,\bar p)\) give the same condition on
\(z\): the two formulas differ only in the name of the free variable being
interpreted, \(y\) occurs in \(\varphi(y,\bar p)\) only at the places where
\(x\) occurred free in \(\varphi(x,\bar p)\), and the remaining variables,
including the parameters, are interpreted alike. Hence
\(z\in b_2\leftrightarrow(z\in A\land\varphi(z,\bar p))
\leftrightarrow z\in b_0\). As \(z\) was arbitrary, \(b_2=b_0\) by
extensionality.
\end{proof}
\end{annexitem}

\begin{annexitem}{annex:p01-w02-p01-005-membership}{thm:p01-w02-p01-005-membership}{Membership conditions of the operations}
\begin{proof}
The pairing and union axioms, with their witnesses made unique by
extensionality, give for all sets \(a,b,c,z\)
\[
  z\in\{a,b\}\leftrightarrow(z=a\lor z=b),\qquad
  z\in\textstyle\bigcup c\leftrightarrow\exists y\,(y\in c\land z\in y).
\]

\emph{Clause 1.} By definition \(A\cup B=\bigcup\{A,B\}\), so
\(z\in A\cup B\) iff there is \(y\) with \(y\in\{A,B\}\) and \(z\in y\), that
is, with \((y=A\lor y=B)\land z\in y\). If such a \(y\) exists and \(y=A\),
then \(z\in A\) by substitutivity of identity; if \(y=B\), then \(z\in B\).
Conversely, if \(z\in A\), the set \(y=A\) satisfies
\((y=A\lor y=B)\land z\in y\); if \(z\in B\), the set \(y=B\) does.

\emph{Clauses 2 and 3.} The intersection \(A\cap B\) is the separated set
\(\{x\in A:x\in B\}\), formed with the \(\Lang_{\in}\)-formula \(x\in B\)
(parameter \(B\)), in which the formed variable is not free. By
\cref{thm:p01-w02-p01-005-separated-set} it exists, is unique and satisfies
\(z\in A\cap B\leftrightarrow(z\in A\land z\in B)\) for every \(z\). The relative
complement \(A\setminus B\) is the separated set \(\{x\in A:x\notin B\}\),
formed with the \(\Lang_{\in}\)-formula \(x\notin B\), that is
\(\neg\,x\in B\) (parameter \(B\)), in which the formed variable is not free;
\cref{thm:p01-w02b-p01-004-relative-complement} writes out this application of
\cref{thm:p01-w02-p01-005-separated-set} and gives
\(z\in A\setminus B\leftrightarrow(z\in A\land z\notin B)\) for every \(z\).

\emph{Clause 4.} Since \(F\neq\varnothing\), \(F\) has a member: otherwise
\(F\) and \(\varnothing\) would both have no members and would be equal by
extensionality. Fix \(Y_0\in F\). The set \(\bigcap F\) is the separated set
\(\{x\in\bigcup F:\psi(x,F)\}\) with
\(\psi(x,F)\defeq\forall Y\,(Y\in F\to x\in Y)\), an \(\Lang_{\in}\)-formula
with parameter \(F\) in which the formed variable is not free, so by
\cref{thm:p01-w02-p01-005-separated-set}
\(z\in\bigcap F\leftrightarrow(z\in\bigcup F\land\psi(z,F))\). From left to
right, \(\psi(z,F)\) follows. From right to left, assume \(\psi(z,F)\). Then
\(z\in Y_0\) because \(Y_0\in F\), so \(\exists y\,(y\in F\land z\in y)\) holds
with \(y=Y_0\), and \(z\in\bigcup F\) by the union condition. Both conjuncts
hold, so \(z\in\bigcap F\).

\emph{Clause 5.} By clause 1 and the pairing condition,
\(z\in\{a,b\}\cup\{c\}\) iff \((z=a\lor z=b)\lor(z=c\lor z=c)\), which is
equivalent to \(z=a\lor z=b\lor z=c\). Likewise
\(z\in s\cup\{s\}\) iff \(z\in s\lor(z=s\lor z=s)\), which is equivalent to
\(z\in s\lor z=s\).

\emph{Clause 6.} If \(z\in A\), then \(z\in A\lor z\in B\), so
\(z\in A\cup B\) by clause 1; the same holds with \(B\) for \(A\). If
\(z\in A\cap B\) or \(z\in A\setminus B\), then \(z\in A\) by clause 2 or 3.
If \(z\in\bigcap F\) and \(Y\in F\), then \(z\in Y\) by clause 4. Each of
these, for arbitrary \(z\), is the stated inclusion.

\emph{Clause 7.} The set \(\{A,B\}\) is nonempty, because \(A\in\{A,B\}\) by
the pairing condition with \(A=A\). By clause 4,
\(z\in\bigcap\{A,B\}\) iff \(\forall Y\,((Y=A\lor Y=B)\to z\in Y)\).
Instantiating \(Y\) by \(A\) and by \(B\) gives \(z\in A\) and \(z\in B\).
Conversely, if \(z\in A\) and \(z\in B\), and \(Y=A\lor Y=B\), then
\(z\in Y\) by substitutivity of identity. So
\(z\in\bigcap\{A,B\}\leftrightarrow(z\in A\land z\in B)
\leftrightarrow z\in A\cap B\) by clause 2, and the sets are equal by
extensionality.
\end{proof}
\end{annexitem}

\begin{annexitem}{annex:p01-w02-p01-005-product-bound}{thm:p01-w02-p01-005-product-bound}{A bounding set for the Cartesian product}
\begin{proof}
\emph{The bound is a set.} The set \(A\cup B\) exists by pairing and union
(\cref{def:p01-w02-p01-005-operations}), and
\(\mathcal P(\mathcal P(A\cup B))\) exists by two applications of the
power-set axiom. Its members are the sets \(w\) with
\(w\subseteq\mathcal P(A\cup B)\), and the members of
\(\mathcal P(A\cup B)\) are the sets \(v\) with \(v\subseteq A\cup B\).

\emph{Every pair lies in the bound.} Let \(a\in A\) and \(b\in B\). By
\cref{thm:p01-w02-p01-005-membership}(1), \(a\in A\cup B\) and
\(b\in A\cup B\). The members of \(\{a\}=\{a,a\}\) are the sets equal to
\(a\), and each of them lies in \(A\cup B\) by substitutivity of identity, so
\(\{a\}\subseteq A\cup B\) and \(\{a\}\in\mathcal P(A\cup B)\). The members of
\(\{a,b\}\) are the sets equal to \(a\) or to \(b\), and each lies in
\(A\cup B\), so \(\{a,b\}\in\mathcal P(A\cup B)\). The members of
\(\langle a,b\rangle=\{\{a\},\{a,b\}\}\) are the sets equal to \(\{a\}\) or
to \(\{a,b\}\), and each lies in \(\mathcal P(A\cup B)\); hence
\(\langle a,b\rangle\subseteq\mathcal P(A\cup B)\), that is,
\(\langle a,b\rangle\in\mathcal P(\mathcal P(A\cup B))\).

\emph{The condition is a formula.} By extensionality and the pairing
condition, \(w=\{a\}\) is equivalent to
\(\forall v\,(v\in w\leftrightarrow v=a)\), \(w=\{a,b\}\) is equivalent to
\(\forall v\,(v\in w\leftrightarrow(v=a\lor v=b))\), and
\(z=\langle a,b\rangle\) is equivalent to
\(\forall w\,\bigl(w\in z\leftrightarrow(w=\{a\}\lor w=\{a,b\})\bigr)\) with
the two inner equations so replaced. Hence the condition
\[
  \pi(z,A,B)\defeq
  \exists a\,\exists b\,(a\in A\land b\in B\land z=\langle a,b\rangle)
\]
is equivalent to an \(\Lang_{\in}\)-formula with parameters \(A,B\), in which
the variable for the formed set does not occur.

\emph{Formation.} By \cref{thm:p01-w02-p01-005-separated-set} there is a unique
set \(P\) with
\(z\in P\leftrightarrow(z\in\mathcal P(\mathcal P(A\cup B))\land\pi(z,A,B))\)
for every \(z\). If \(\pi(z,A,B)\), then \(z=\langle a,b\rangle\) for some
\(a\in A\), \(b\in B\), and the previous step gives
\(z\in\mathcal P(\mathcal P(A\cup B))\). So the first conjunct follows from
the second, and \(z\in P\leftrightarrow\pi(z,A,B)\). The members of \(P\) are
therefore exactly the ordered pairs \(\langle a,b\rangle\) with \(a\in A\) and
\(b\in B\), which is the membership condition of the product
\(A\times B\) of \cref{sec:set-prerequisites-02}. Any set with that membership
condition equals \(P\) by extensionality.
\end{proof}
\end{annexitem}

\begin{annexitem}{annex:p01-w02-p01-005-laws}{thm:p01-w02-p01-005-laws}{Laws of union, intersection and relative complement}
\begin{proof}
\emph{Reduction.} Fix an arbitrary set \(z\) and write
\(p\), \(q\), \(r\) for the statements \(z\in A\), \(z\in B\), \(z\in C\).
The empty-set axiom gives \(z\notin\varnothing\). By
\cref{thm:p01-w02-p01-005-membership}(1)--(3), applied repeatedly, membership
of \(z\) in each side of each identity is equivalent to a propositional
combination of \(p,q,r\) and the false statement \(z\in\varnothing\). Each of
\(p,q,r\) is either true or false. For each identity below it is shown that
the two combinations have the same truth value for every choice of truth
values of \(p,q,r\) (in clause 10, for every choice compatible with the
hypothesis). Then \(z\) is a member of one side iff it is a member of the
other, and since \(z\) was arbitrary the sides are equal by extensionality.
In each case the argument splits on the truth value of one statement and
evaluates both combinations.

\emph{Clause 1.} The sides of \(A\cup A=A\) give \(p\lor p\) and \(p\); the
sides of \(A\cap A=A\) give \(p\land p\) and \(p\). If \(p\) is true, all four
are true; if \(p\) is false, all four are false.

\emph{Clause 2.} The sides give \(p\lor q\) and \(q\lor p\), and \(p\land q\)
and \(q\land p\). A disjunction is true iff at least one of its two components
is true, and a conjunction iff both are; these conditions do not depend on the
order of the components.

\emph{Clause 3.} The sides give \((p\lor q)\lor r\) and \(p\lor(q\lor r)\):
each is true iff at least one of \(p,q,r\) is true. The intersection sides
give \((p\land q)\land r\) and \(p\land(q\land r)\): each is true iff all three
are true.

\emph{Clause 4.} The first identity gives \(p\land(q\lor r)\) and
\((p\land q)\lor(p\land r)\). If \(p\) is false, both are false. If \(p\) is
true, both have the truth value of \(q\lor r\). The second identity gives
\(p\lor(q\land r)\) and \((p\lor q)\land(p\lor r)\). If \(p\) is true, both are
true. If \(p\) is false, both have the truth value of \(q\land r\).

\emph{Clause 5.} The sides give \(p\lor(p\land q)\) and \(p\), and
\(p\land(p\lor q)\) and \(p\). If \(p\) is true, all four are true; if \(p\) is
false, \(p\land q\) is false, so \(p\lor(p\land q)\) is false, and
\(p\land(p\lor q)\) is false.

\emph{Clause 6.} Write \(e\) for the false statement \(z\in\varnothing\). The
identities give, respectively, \(p\lor e\) and \(p\); \(p\land e\) and \(e\);
\(p\land\neg e\) and \(p\); \(p\land\neg p\) and \(e\). Since \(e\) is false,
\(p\lor e\) and \(p\land\neg e\) have the truth value of \(p\), and
\(p\land e\) is false like \(e\). The statement \(p\land\neg p\) is false
whether \(p\) is true or false, like \(e\).

\emph{Clause 7.} The first identity gives \(r\land\neg(p\lor q)\) and
\((r\land\neg p)\land(r\land\neg q)\). If \(r\) is false, both are false. If
\(r\) is true, the first is true iff \(p\lor q\) is false, that is, iff both
\(p\) and \(q\) are false, and the second is true iff \(\neg p\) and
\(\neg q\) are both true, the same condition. The second identity gives
\(r\land\neg(p\land q)\) and \((r\land\neg p)\lor(r\land\neg q)\). If \(r\) is
false, both are false. If \(r\) is true, the first is true iff \(p\land q\) is
false, that is, iff at least one of \(p,q\) is false, and the second is true
iff at least one of \(\neg p,\neg q\) is true, the same condition.

\emph{Clause 8.} The sides give \(p\land\neg q\) and \(p\land\neg(p\land q)\).
If \(p\) is false, both are false. If \(p\) is true, \(p\land q\) has the truth
value of \(q\), so both have the truth value of \(\neg q\).

\emph{Clause 9.} By the definition of inclusion, \(A\subseteq B\) states
\(\forall z\,(p\to q)\). By extensionality and
\cref{thm:p01-w02-p01-005-membership}, \(A\cup B=B\) states
\(\forall z\,((p\lor q)\leftrightarrow q)\), \(A\cap B=A\) states
\(\forall z\,((p\land q)\leftrightarrow p)\), and \(A\setminus B=\varnothing\)
states \(\forall z\,((p\land\neg q)\leftrightarrow e)\), which, as \(e\) is
false, is \(\forall z\,\neg(p\land\neg q)\). It suffices to show that for each
\(z\) the four matrices \(p\to q\), \((p\lor q)\leftrightarrow q\),
\((p\land q)\leftrightarrow p\) and \(\neg(p\land\neg q)\) have the same truth
value. If \(p\) is false, \(p\to q\) is true; \(p\lor q\) has the value of
\(q\), so the second is true; \(p\land q\) is false like \(p\), so the third is
true; \(p\land\neg q\) is false, so the fourth is true. If \(p\) is true,
\(p\to q\) has the value of \(q\); \(p\lor q\) is true, so the second has the
value of \(q\); \(p\land q\) has the value of \(q\) and \(p\) is true, so the
third has the value of \(q\); \(p\land\neg q\) has the value of \(\neg q\), so
the fourth has the value of \(q\). Equivalent matrices for every \(z\) give
equivalent universal statements.

\emph{Clause 10.} The hypothesis \(A\subseteq C\) says that \(p\to r\) is true
for every \(z\); so the combination in which \(p\) is true and \(r\) is false
does not occur. The identity \(C\setminus(C\setminus A)=A\) gives
\(r\land\neg(r\land\neg p)\) and \(p\). If \(r\) is false, the first is false,
and \(p\) is false by the hypothesis. If \(r\) is true, \(r\land\neg p\) has the
value of \(\neg p\), so the first has the value of \(p\). The identity
\(A\cup(C\setminus A)=C\) gives \(p\lor(r\land\neg p)\) and \(r\). If \(r\) is
false, \(p\) is false by the hypothesis and the first is false. If \(r\) is
true, the first is \(p\lor\neg p\), which is true. The identity
\(A\cap(C\setminus A)=\varnothing\) gives \(p\land(r\land\neg p)\) and \(e\);
the first contains both \(p\) and \(\neg p\) as conjuncts and is false, like
\(e\). Finally, \(A\cup C=C\) and \(A\cap C=A\) follow from
\(A\subseteq C\) by clause 9 with \(C\) in place of \(B\).
\end{proof}
\end{annexitem}

\begin{annexitem}{annex:p01-w02-p01-005-subset-algebra}{thm:p01-w02-p01-005-subset-algebra}{The subsets of a set form a Boolean algebra}
\begin{proof}
\emph{Carrier and closure.} \(\mathcal P(S)\) is a set by the power-set axiom,
and \(X\in\mathcal P(S)\) iff \(X\subseteq S\). Let \(X,Y\subseteq S\). If
\(z\in X\cup Y\), then \(z\in X\) or \(z\in Y\) by
\cref{thm:p01-w02-p01-005-membership}(1), and in either case \(z\in S\); so
\(X\cup Y\subseteq S\). If \(z\in X\cap Y\), then \(z\in X\) and so
\(z\in S\); so \(X\cap Y\subseteq S\). By
\cref{thm:p01-w02-p01-005-membership}(6), \(S\setminus X\subseteq S\).
The empty-set axiom gives \(\varnothing\subseteq S\), since no \(z\) satisfies
\(z\in\varnothing\), and \(S\subseteq S\) holds because \(z\in S\to z\in S\).
Hence \(X\cup Y\), \(X\cap Y\), \(S\setminus X\), \(\varnothing\) and \(S\) are
members of \(\mathcal P(S)\).

\emph{The operations are functions.} Write \(Q=\mathcal P(S)\). By
\cref{thm:p01-w02-p01-005-product-bound}, \(Q\times Q\) and
\((Q\times Q)\times Q\) are sets. The condition
\(\exists X\,\exists Y\,(X\in Q\land Y\in Q\land
w=\langle\langle X,Y\rangle,X\cup Y\rangle)\) is equivalent to an
\(\Lang_{\in}\)-formula with parameter \(Q\), after the ordered pairs are
unfolded as in the proof of \cref{thm:p01-w02-p01-005-product-bound} and
\(u=X\cup Y\) is unfolded as
\(\forall v\,(v\in u\leftrightarrow(v\in X\lor v\in Y))\), which is
equivalent to it by \cref{thm:p01-w02-p01-005-membership}(1) and
extensionality. By \cref{thm:p01-w02-p01-005-separated-set} the members of
\((Q\times Q)\times Q\) satisfying it form a set \(G_{\cup}\), and every triple
\(\langle\langle X,Y\rangle,X\cup Y\rangle\) with \(X,Y\in Q\) lies in that
bound because \(X\cup Y\in Q\). So \(G_{\cup}\subseteq(Q\times Q)\times Q\) is a
binary relation from \(Q\times Q\) to \(Q\). For each member of \(Q\times Q\),
which is \(\langle X,Y\rangle\) with \(X,Y\in Q\), the value \(X\cup Y\)
satisfies \(\langle\langle X,Y\rangle,X\cup Y\rangle\in G_{\cup}\). If also
\(\langle\langle X,Y\rangle,u\rangle\in G_{\cup}\), then
\(\langle\langle X,Y\rangle,u\rangle=\langle\langle X',Y'\rangle,X'\cup Y'\rangle\)
for some \(X',Y'\in Q\); the ordered-pair identity of
\cref{sec:set-prerequisites-02}, applied twice, gives \(X=X'\), \(Y=Y'\) and
\(u=X'\cup Y'=X\cup Y\). So \(G_{\cup}\) is a function from \(Q\times Q\) to
\(Q\). For intersection, \(G_{\cap}\) is the separated set
\(G^{Q}_{\cap}=\{w\in(Q\times Q)\times Q:\gamma_{\cap}(w,Q)\}\) of
\cref{def:p01-w02b-p01-004-graph}, formed with the \(\Lang_{\in}\)-formula
\(\gamma_{\cap}\) of \cref{def:p01-w02b-p01-004-formulas} with parameter
\(Q\); by \cref{thm:p01-w02b-p01-004-intersection-on-subsets}, whose proof
verifies the bound through \(X\cap Y\in Q\) and totality and
single-valuedness through the ordered-pair identity, it is a function from
\(Q\times Q\) to \(Q\) with value \(X\cap Y\) at \(\langle X,Y\rangle\). For
complement, the members \(w\) of \(Q\times Q\) with
\(\exists X\,(X\in Q\land w=\langle X,S\setminus X\rangle)\) form a set
\(G_{c}\); for each \(X\in Q\) the value \(S\setminus X\in Q\) belongs to it,
and if \(\langle X,u\rangle=\langle X',S\setminus X'\rangle\) with \(X'\in Q\),
then \(X=X'\) and \(u=S\setminus X\). So \(G_{c}\) is a function from \(Q\) to
\(Q\).

\emph{The laws.} Take meet \(\cap\), join \(\cup\), complement
\(X\mapsto S\setminus X\), bottom \(\varnothing\) and top \(S\). For
\(X,Y,Z\in Q\), \cref{thm:p01-w02-p01-005-laws} gives: commutativity of both
operations (clause 2); associativity of both (clause 3); the two absorption
laws (clause 5), and idempotence (clause 1); both distributive laws
(clause 4); the bounds \(X\cup\varnothing=X\), \(X\cap\varnothing=\varnothing\)
(clause 6) and \(X\cap S=X\), \(X\cup S=S\) (clause 10 with \(S\) in place of
\(C\), using \(X\subseteq S\)); and the complement laws
\(X\cap(S\setminus X)=\varnothing\) and \(X\cup(S\setminus X)=S\) (clause 10).
These are the bounded distributive lattice and complement laws.

\emph{Order and nontriviality.} By the definition in
\cref{def:p01-bpi-boolean-algebra}, \(X\leq Y\) iff \(X\cap Y=X\), and by
\cref{thm:p01-w02-p01-005-laws}(9) this holds iff \(X\subseteq Y\). The algebra
is nontrivial iff its bottom and top differ, that is, iff
\(\varnothing\neq S\).
\end{proof}
\end{annexitem}

\begin{annexitem}{annex:p01-w02-p01-005-no-universal}{thm:p01-w02-p01-005-no-universal}{Collections with no bounding set}
\begin{proof}
\emph{Clause 1.} The formula \(x\notin x\) is an \(\Lang_{\in}\)-formula with no
parameters, and the formed variable does not occur in it, so
\(r_a=\{x\in a:x\notin x\}\) exists by
\cref{thm:p01-w02-p01-005-separated-set} and satisfies
\(\forall x\,(x\in r_a\leftrightarrow(x\in a\land x\notin x))\). Suppose
\(r_a\in a\). Instantiating \(x\) by \(r_a\) gives
\(r_a\in r_a\leftrightarrow(r_a\in a\land r_a\notin r_a)\), and since
\(r_a\in a\), \(r_a\in r_a\leftrightarrow r_a\notin r_a\). If
\(r_a\in r_a\), the biconditional gives \(r_a\notin r_a\); if
\(r_a\notin r_a\), it gives \(r_a\in r_a\). Either case is contradictory, so
\(r_a\notin a\).

\emph{Clause 2.} Suppose \(\forall x\,(x\in V)\). Then \(r_V\in V\), which
contradicts clause 1 applied to \(a=V\).

\emph{Clause 3.} Suppose \(\forall x\,(x\in D\leftrightarrow x\notin A)\). Let
\(x\) be arbitrary. If \(x\in A\), then \(x\in A\cup D\). If \(x\notin A\),
then \(x\in D\) and again \(x\in A\cup D\), by
\cref{thm:p01-w02-p01-005-membership}(1). So \(\forall x\,(x\in A\cup D)\),
contradicting clause 2 with \(V=A\cup D\). The class \(\{x:x\notin A\}\) is
therefore not a set.

\emph{Clause 4.} Suppose
\(\forall x\,(x\in I\leftrightarrow\forall Y\,(Y\in\varnothing\to x\in Y))\).
For every \(x\) and every \(Y\), \(Y\in\varnothing\) is false by the empty-set
axiom, so \(Y\in\varnothing\to x\in Y\) is true; hence
\(\forall Y\,(Y\in\varnothing\to x\in Y)\) holds for every \(x\), and
\(\forall x\,(x\in I)\), contradicting clause 2. For the second sentence:
\(z\in\bigcup\varnothing\) iff \(\exists y\,(y\in\varnothing\land z\in y)\),
which is false for every \(z\); so \(\bigcup\varnothing\) has no members. The
set \(\{x\in\bigcup\varnothing:\forall Y\,(Y\in\varnothing\to x\in Y)\}\) is
included in \(\bigcup\varnothing\) by
\cref{thm:p01-w02-p01-005-separated-set}, so it has no members either, and it
equals \(\varnothing\) by extensionality.

\emph{Clause 5.} Pairing applied to \(\varnothing\) and \(\varnothing\) gives
\(a=\{\varnothing\}\), with \(z\in a\leftrightarrow z=\varnothing\); in
particular \(\varnothing\in a\). Suppose some \(b\) satisfies
\(\forall x\,(x\in b\leftrightarrow(x\in a\land x\notin b))\). Instantiating
\(x\) by \(\varnothing\) and using \(\varnothing\in a\) gives
\(\varnothing\in b\leftrightarrow\varnothing\notin b\), which is contradictory
as in clause 1. So no such \(b\) exists for this \(a\), and the universal
sentence is false.
\end{proof}
\end{annexitem}
